% Standalone paper: mathematics, references, data and verifier.
% Compile three times with pdflatex -no-shell-escape.
\documentclass[11pt]{article}
\usepackage[T1]{fontenc}
\usepackage{lmodern}
\usepackage[margin=1in]{geometry}
\usepackage{amsmath,amssymb,amsthm,mathtools}
\usepackage{booktabs,longtable,array}
\usepackage{xcolor}
\usepackage{microtype}
\usepackage{listings}
\usepackage[colorlinks=true,linkcolor=blue!55!black,citecolor=blue!55!black,urlcolor=blue!55!black]{hyperref}
\usepackage{attachfile2}
\hypersetup{pdftitle={A simpler construction of square-difference-free sets beyond exponent 0.758},pdfauthor={Eric Naslund}}
\newtheorem{theorem}{Theorem}[section]
\newtheorem{lemma}[theorem]{Lemma}
\newtheorem{proposition}[theorem]{Proposition}
\newtheorem{corollary}[theorem]{Corollary}
\theoremstyle{definition}
\newtheorem{definition}[theorem]{Definition}
\theoremstyle{remark}
\newtheorem{remark}[theorem]{Remark}
\newcommand{\Z}{\mathbb Z}
\newcommand{\N}{\mathbb N}
\newcommand{\R}{\mathbb R}
\newcommand{\QR}{\mathcal Q}
\newcommand{\al}{\alpha_0}
\newcommand{\card}[1]{\lvert #1\rvert}
\DeclareMathOperator{\supp}{supp}
\lstset{basicstyle=\ttfamily\small,columns=fullflexible,breaklines=true,keepspaces=true,showstringspaces=false,frame=single}
\begin{filecontents*}[overwrite]{sarkozy-simple-certificate.py}
#!/usr/bin/env python3
"""Complete finite certificate for square-difference-free exponent 0.758001.
Run with Python 3 without -O: standard library only, no network or optimizer.
Default: extract to a temporary directory and execute all checks.
Use --extract DIRECTORY to retain the unpacked data, verifiers, and report.
"""
import argparse, base64, hashlib, io, pathlib, subprocess, sys, tempfile, zipfile
ARCHIVE_SHA256 = '4e768df0fcf2096bdde3416928a709a6ba1a6119231bc27510a1e377a50140c3'
PAYLOAD = """
UEsDBBQAAAAIAAAAK11SvVPYQwEAAHgCAAANAAAAbWFuaWZlc3QuanNvbn3Q22obQQwA0Hd/xbLP
wdXMSHPpr4RSpJFEXFyvawdCCfn3zm4CSSDuixh0m4Oed9M08/H8wPP3aYZ9oQoQ5rs1vaj+7Mvv
83Ky0+N11O9HdpqetzjqfjjaOqX8yN/+xED7X9fltM1u9esDR8prhxgFKKkXAAPXopa7mKA1447m
7I5kGANbcoitBtGCATLmzKX5+87z8mSXdeUeYq7zln65+58LU7npihzMS5KWigQD1VwHASkVKK2h
Onk1KYCu1ASqddXaumZPKKnLl65c0ptrxB/bJeVw4svf92OOxlftZ+tr20ftB6tTyB4KUSXQSFYq
sibsHUIVGE/1DGwsWTNjai0ReApIUaMJ+GraLjUfD3IZ/9xAvFVvKFRBe8mayFoSSa2XANxCiK2p
ICGqMltH4k7D5cGpaAvOVSjWuCl2L7t/UEsDBBQAAAAIAAAAK13cWhuDvAEAABQHAAAOAAAAZGF0
YS9xMjE1Lmpzb26tlc1uozAQx+95Csun7ipp8Bfg3vcpVhWaYFO8BZs1zqZS1XdfA00TbZU0kpcD
nvF/+OEZD+Z1hRAevOn1iB/Qz+ghJNbzwFkcHicbKz2ENupk9gZnbDiFL/fTiFD2bj2uP0mkPGqr
s4jPCHYZQcVtCHIFUdyG4FcQ8jbElVpQno6gyeXkNyLoZUT6hhCZvAiepXcWTW4LJpJXwZK7grH0
PGT6R1Ym7yn9Z08/DqSDUaGtlLauNxaC89PhlMXrTLX7XvtJOx1VhJRk/R8tUcgvDcKLfLFkyS9O
ReMjOf0Cdah6eKkGOM8jpiGlpHNMAP+kQwXd0EKcx/eFKLOM4EXUY9CqGtxB+1nMaF4u0hgg7KeC
4B/TW1Dd6vpZKzSafuhMY2oIxtkHpNzBHsAr5N3eKmOfkGvQCH9iaKNhNLtOI6cUmks9vrPd3tc6
sl/nlHBjusnDXo8afN1urIsr28bHNsvE9jclYjP/XMj9r9FZvBQDjy1QkU/PFjQjIBUnmVA6JzvG
CM+EZEoJqmQGrBBUcAmMSKKaBgA431GVs1KLmkBzJB7TmJhN55y/m9vl+7z+b9vZwTH0bfW2+gtQ
SwMEFAAAAAgAAAArXXY/Ht+JDwAAGRwBAA4AAABkYXRhL3E0MzcuanNvbsVd23LbOBJ9z1e49LS7
NZlBN3HN+37F1lRKY8lj7yaW13Y2UzU1/77yBZDZtJQGiEPlIaBA8vjwsNFoAC3ozw8XF6u7+5uv
24fVp4t/7T9dXFD66bnkYV/8+nS82mzvHq/3F/Dzp7vdze3j4fqX/w/lHsH89Ob49fDXnybX+cNl
Pl/14c21P0C2KGSGUdaJQakeWonsYcgNnInOTJoNjjROjybWOkFaSHudHgYHnWB6MDdAh3MLgmSd
YI46wUg36WFh0IySg3FqKFsiozy1a+AMo6z00w7XCrWkbX3sYRqgLQyacaT7Gp6tDphgyA3A2qAG
p0aAkQbKEWBm12LQHuc7PExppXnEBuhwbj2ArB3MSxuc0hHXtURU14LrsyLO7gaUm2ZqIA3jrB0v
E86klay5fuRp6pEtCplhlLtGNFwddTAIuJ4w0XkZa90Ro+KvBi0iyt5amrXONcOAHcwRUT1yOK8W
OMoO5ZRbLC7AxLAoZEaJwTAtfM/FoGq3HOoJg/gqnXJA2XHn5ba8yglYfJwJzTjSSj06rz+OoD0O
uvOK20K0e69BLiVJ71XIt9idlyFHkhggdsJJ0ntNbylNoLwTznEnHO3ey5FzsRmmCAMF0bZJhnnu
vmuSc0lr/bYDtkclba37q4/+tC1mFjLDKCvF0Po9pUk3pFo1INdzPmnRi5DWxjgtpHF6tLBWCtJA
2iv1MDjoBNODGzxeOLsgSNYJ5qgTjHSLHhYGzSg5GKeGtiUyylO7es4wylo/7XCtsK/O6b1nUyMn
GDLBkPtOHJp37ugz31SPXA8czk1Zu85XzxknRoQ1wGpg5cyBdtGlHtnDtOgcQy+iBpC0h/lmnM4J
xjnB+hMYcIIZnUU555ZocYCRDjjS1Jc0EW7eZ4RtcdgMpK2UpEFtrdgEg24xETo7a3W2ewNtpCIB
2CC5AVs7SYODBroo9aowNUzTAMUG0rZAp81AtQ2Odws24yRhoCK+61xQg9v2DaRhnLVO2+Ps2vUc
jbIBjb7mARMKWClFw0DGK5E9DDn1HG0sw7lhpS+cW40Gzkqr6zp6ZjNrKc7CZFaK0XmIu4wcSNIJ
5Z4TjHODHAknNEwNGGVtI2SUf3Zd58Zmcg4wztSXs3rcmeYFjBYGzTjSvZPhCZfCT7gUfgKm8BMu
hZ+AKfwETOEnXAo/AVP4CZjCT8AUfgKm8BMwhZ+AKfyES+EnYAo/AVP4CZbCT7gUfgKm8BMshZ9w
KfyES+EnXAq/1v+1RPDaRjMXm4G0tZJ0teyW8W4DdGfLXoi1xbEGCmJxtuf6emw21V8Cb4EGNsbY
160uJAiQNdBfOxzriGMdYZ2MgyFHnOl1TjZpWSVQRmZzWat7AoOzazVv5feYcBtgwPa/wG1/gdv9
AreLBHCjh0VI9975Yhk9eu97gdv2ArjDA3DTC+DeEcsIgmSdYI46wUj33u4Ct9sFbLML4F4XuK0u
YDtdADe6wO1z0ULaK9cI6lOyRtAWBs040ko91B7P6+Z+G9KyWqAbWJ+c/F2Gtjq0aaENlKSJt1KT
FtpeKYkBYiecJNzi/8L5NYHyTjjHnXC0mySxOGyGKcJAQbRtkmGe2zWwxpHW+m0HbI+9Y7TU8HsC
vmHxn/qa9lzavVNbPC65qgG6t20vQ5sckDZQkibeSk16hw1pXsDdPbWqQRLq7F0XkgRIu3u0M9MA
A04RC4NmmB6Mk0PbHGE+u3OANo9z7zyzlqbYd9GAGbaYOROacaSVenRez2SGrQ3OZX1ylX4Z2r3X
NJeSpPeq5lvszsuazLhlvNlya82k8xrhUppAeSec40442r2XN+diM0wRBgqibZMM89x91zjnktb6
bQdsj1raAbiEEIDzzgE3O9wiST1079XIeugWPejsrC2ONVAQi7M91wCtXPQIOGhgY4wN0OHsggBZ
A/21w7GOONYR1sk4GHLEmZ46OT3gslJC/eImAXsCg7Nr6ryU1xS+e+B4xuNGHQ2StMhtOg/WG6Cp
78zwUrQTkDZSkoQzwM6rm2+hIw4a6aWocy7XUpIAaRPhHLcB0gb2N4TrbwwMmRhofzi3bbuucM4k
rXbaFmjYnYenad7cCFkcNgNp987lCrgNzwIuAy0ANzxrUATHGqoI0Pw6D63TvHk/DxS7c8ZLqE8e
ASoCZN19zWOm+QWcIhYGzTA9GCeH0vYizl+HvlNo80ir3XXAWbWWtoVtamBV+xTPRO69p0G1HEM9
cNJ9CdeikBso05kpBxRlmBYBZXCpHlj57foBhTygtFB/4cKqdglfRAwc5wHlmBNMZYei3IKs7EsS
DLmFtHLLFoJ1U75BDgMjTTDSoS/noPxJhLn5ahaITUBspSaua9zftNxVD91COp6ddcSxBgoScbbX
0tYdbGrd4aQeOq81hIblUJzUQNYO57CZgGJbYEdjcR0NE7ATszALDDCX3dAZeBznCCOtNey+g/KW
rzfqBjIzkbU27WFyUOo63mj6bmMDdOo64FiIddO06dkF6Twj2+JAbLXXI4ODTjA91OtyutHuQoIg
WSeYt04w0i16WBw0EUwQdbjHuNbIMHft+s6XzWVNMNad0yO5evPkoR7YoYAZBdx3Opxrd3vWDovq
gesZm/MyNijCKCUMytrqFe67/zej1uQaGp6vR6Yzi4Gj7FAe2cIoexTlFmRGicEoyjqTg/VOsRo4
gAgbFOG+vxbADVFn335kJjTjSHee++L6WTUYcOpqdotwZhxlnBoMs7qG9t35W46M24GMcWtYjPse
YosgQNYJ5qYdjrTF9S0W1bc4WKdlYYZnUG66gTKKsdZJ49pg759xaVjL0+6hPA+acaSVevT+XaWW
tbwG6N4/IrQM7e6/q7SQJAxskL1/V2nm5iRq7ISTpPvvEy2kCZR3wjnuhKPd/ZcGcJsozVSEgYJo
2yTDPHfnbfvn7vJBONahL23lntiuOthR7qI8C5hRhHVKaH2ebo9jVx/i1APXMz5ly0tQ1sY2DZRh
YrRw1qnRQNnrxDAw5IQSgxt8XDizGkDOCeWYE4pyixgWhcwgLRgmhbL9Mcgzu3rGKMJKv+xgbU8b
JDZMALl6aAuDZhzp3rNW9ZNWFoZcz1k9rQQjnXCccXIknN2Feuje82AN0B4nSKyHDucXBMja4zx1
wLGOONYR1r8EGHLEmV5CuWpuCD9gnNWLFoyzaS1rgs35EFWPWVqgGUdap4etR1YKTSjkBsraGSUY
Z7XDqycNlCPgmiE3QCtnaGDIFiaHegGY6udocELjSFuck2ac0gbGugWaYXowTg7fcx6o3k37Bsoo
xkon7WH23LJrTOcssYDLEguwLLEAzBILuLmrgMsSC8A5poDLEgvALLEAzBILuCyxAJwdC8AheQBm
iQXgBFkAZokF4MA84LLEAjBLLACzxAIsSyzgssQCMEsswLLEAi5LLOCyxML8LLEPr3esvt9sHq8/
b7a3u683t+vH3f3q097p7P+9OXv77ev2/uncw/7kCz4PL9hDGJfhNSz0uV58tn5c5i8bRPE5n4/i
+oyXS/uD0r3yTDT+bF+XfmMuX6+P6Uj5eh2Z1wOXREV+8iAYkqVxRZAncoUL4zIMp8v8jORiFs/8
sCbzTnZcUny9gIIXT5grkhmXlgV2MYJ8J4mKyRPa06XNnPJBLgO//zm/RiJ5wENWorAzSdYMXhz4
VyLsREW5opw53GysPBgUBzn3k30SNak8Qj4wNG5VZJw44DyMK+81RfmiiyL5wL/ew9kyc0W5opw5
3GxyC8uP78sBTw7yqcxuYCdqOLfMaU2WKrfhXJIRQjkz/uxFOeSHyN4mV+QLygkvEHOZzWBSOgmU
xi2g6OeFg+Pc3KYH1krZpq/unZfJQr7yFg4vU96T0rgs9A8H+SdbqNxsC5w8FbJ5OmkJ+akTj0vK
vu2gjztdFhpBenYnTuSKjKz9C64A0bglRfv+54O9Gy8PyjVOKHGsLAaT78wVUULmCh/HZQynS18U
tOOKkB0Fe9HXlQPpl8vj5VtKyzJj9uWxkjAV0U5yfelEeVxKPFlmHpHHn4vRys/ZtZIz4mDID83F
8ZS/Umry3ykH+a0MQ5LvjcWZw82lZph0UacOSowku/LSCw6hNNYsSBZ2GJdeWAf8Rf4gkMvOPwcf
uTGWXtfR+3FYEB6YSlxT7ohjJ1QalBUuxouKQ4MXd8huQCLKsryw0maTiBVLp+1IvMCpyQ2HeMzo
zSc6cffBgaaJXaU5hpokz+xAOYdxk4pjASLZEhIN49CyVEyi1kF4aSsuLCdyRQ7IcpnfxbGykAsi
zC3BY44ZU5CmKMpJ0ObEOEU6+6HE5GFckdtz7trLgCyNSz+cLh3JTqEg5L+Y91QskWcYB+m5pMMA
xosr7PsjnsOYSgyySthgRUX1mCq/kTJQDeM3MvlMRyK/EpDlML4E07kiijHoEIWiuaIMPq0YjXrp
Pw5DD5oc2PG9pYFQ9n+locYcYpBwVyJyKBEDcRzfUSpkRGlL83TjihJyenGHbPJ5nHGsPIzDSlA5
iAHstKIMQ6YHyQlVgujZZQeTR8Hlr0yGxYO404kyd9vHyjLsc4McB9KRimk4czgoQTkdC/PeCfxK
Nx9FRbk0yZriDsP7tiHLMrIutlGehifmXRyiPTK/kRtTdmDOnS6LfxtERRkPDOIOa8el9D2yLIZy
iM7GDy4/E7HoWcrB4Fh43+k104PcsMvd0R7DPdxc3FQ5UPyl4vQGJyc8StAzrRmEtRVnISO0I/FU
mQYrpigqDo5T3FHc/3C6R8r1pTHlG0tFQZYVGVp24TRIzbyULA7jskRfxXLS+IJyIlfIiYvs5Y6V
VELJPJNaakqMMqko73NywEZMqMg31uENoltjiqdLLnPhg6iYzhDmjp2zQDm6n0zW5Qr5Ast480hZ
RqhezKLleODwuawNbP9YXz5+/rr+4/Pd+u0ywOrT3ke/dIOrx/X979vHz+svd9frff3q5+CiMbR6
Obl9eNxuPt/tvm/vn08aH4aXUw+P68dvT+sJq38+/ZWLy+vt5X+2m4uHm693X26ubi7Xjze7208X
m9332+/r+83F/e7b7ebm9veL3dXFw/p/+0uvtuuHm9++bC92m83F80rFwyv27tv95XaP/efzE62u
br48fVrdbx+26/vL64+3uz2zX/a3fXyp+OW/dggfN9u7x2v+uL66urndftz/t/7y878fdrerF2FW
D9drdv4JyFzx1m5jcFcbE4ar9Xq4cjFyIAoUtsNvZuvNpb1cp43d8voquM1VMBv322VItDXWZcT8
TE+YV192u/u/PS+9/OP5Yf7+y/OH1f7Svz789eH/UEsDBBQAAAAIAAAAK13hpb1HXAIAAMUcAAAQ
AAAAZGF0YS9iaW5hcnkuanNvbu1ZXY+qMBB991cQnr06bSkf9/6UzcZ0sVyasGBq1Riz/32LClvZ
dQMCKkYeTOdQxp6Z6WEadiPLspeKqdXS/mvZkUiF4pZkSmQpS6w3kTK5tUIulYhEyBT/Z8lVegS2
k8XWUpm15lIb9jj3xZJFzHJXMPGoD4AOcHSAEHWO03IbT3wX8AGY84WKNYggv/aQzDKVI3tjzUOV
SW2+aEvbbhBgND3M3nvQGCoGGFyCppga9xwPKJlSc7pDMFSggJAgqHj1PYRIgWnodVzEjC/L9ez2
vxrfiHSebb5u5BeMyyH6Gv45M0bH0WsB2W+SpWHMT53uylEeKn3L+BeNhLFI5kX0SnQZsoSfuKks
ylzAySIOj8ciUt8eh9qPb9ji20IljxKd2xwv4Y/xeZr4gWgWmR4ZpC+opJ/RbusIagTYCSoh8vy7
LiX0I1NSgymtLNQPeiaKzhBFLfYMvQZREtR30EVOycXVi52WVF3aA9U2EoGNvFNTGQBOjKsLBapI
MYUB6sRFighNmVLUwMPvWtHm/fqAWSXDymrf7QTUE45OtQI/WnOKOts8zWuqyfbrrxF37y+l6EYy
0T6luJLRG2sEMWLhwB31E4S0jfRQZILUrohOCupWrcTQEkqGktC+TxyG/JvyZx4j6fO8cZE8oObV
hIfZQ1yDKfH66MD76SMqmxw3FQlE6ju4/mHjl66iU6FwnoeN4vLavl0bvZ/7Ewp4HjbuPKWnKlF+
PYtEypLZRsxVPCv42GhafAsUyzBbc7md/ZfZRk9JV+9cijCvg/yTIfEBfJc6KADPHX2MPgFQSwME
FAAAAAgAAAArXT0lBSWsAQAAhAoAABEAAABkYXRhL2xpYnJhcnkuanNvbs1VzW6CQBC++xSEM2l2
Zv99lcYYU2nLQW3ApAfjuxdlZlkDWcHYphcZh2/m+8HF0yLL8qYst2sEnS+z1/Z7loEXxbUwrruC
Nl2BnhrACE0NZKi1VHigGUDqCMYY1RWSEIK2SFpvgXdge10VQaWSNqhkLZrWKxQskteZsIcawHcU
D0EYMuyExKGRbBZYHoGldaycg0EeV+RfIjAF7eEhqdg+qxHU0LxWh+BFCOBYb/bN+6HeNTcxdJ8t
tKBCFIPO9boq7uGRKztxQIYKJk7MlZTEW75rhhbcbAtTKZ5qYYjHtKDfznQE31u3/8SyC5NzGdIO
PBcqKPJjA3rgVE7MyE/LSA+EB8sTT5qZG9Ef/ozgOWfzvqbw7vyqq125fvvcVPsmvDlPBO7utu3A
l2+rj+rYI29I+BEEPZcNh++ybuH5C7jLP+m1fy7GaexUmr7EBKNwRsg0I8B8StWXOuVXWGXS7PIB
dhexRzH0pY/GRDIdr0Ran/YPPJBIINioj6MJoonqaCP6lHSFCu88WCEf0B5l53E0XQPj9mS00csI
E8eR9CTQ8PG4nM7FefEDUEsDBBQAAAAIAAAAK11tQ+nkSQkAAJwZAAAJAAAAdmVyaWZ5LnB5lVlb
b9s4Fn73r+B0sZDUKortpE3jaQYoOg1QYGZSpJ2nbCDIEm2z1W1IqrFb9L/vd0jKkiy32c2DIfFc
eO4X5V+/nDZKni5FecrLL6ze6U1Vnk2ePHnypirqnGvO+DZJNVuJUuDlC5diJdJEi6pkq0oyveFM
CUKVwKyrkpeaTaOL5y+n01k0mfytuGJVme/Ye8PbU0zppMwSmbFcLGUidwyvrKxYyfVDJT9H7Hee
iiLJWZrkaZObu9SkllVdKc7q6oGuKtO8Uo3k6leW5DlL0pTXmmdMlPyfJsmFFri3UXvxJX6JTzT5
CIlF6fRJcaRlY2BGjlystCjXLOdFkSiWSM6KBEriB2rjIrluCuiowGNCyuNmWCopd0RVJzWXIZTR
LFGKS3Ml49B4mQu1gXzLHUwmFIM6a5kUEZl6AvtVUrNNojawSfv6SVXlZCWrgmXOIA7g7BOy25u/
//o9vv7j5uY2ZHlF9qpKzbfaku2Vbgmv3QGEY9cWh5RrwUL9Ix1pjWOI0kLe43XSypWSXqtd+2pC
YhensJ7IEs0tg3TD08/xmlcF13Cxw1VF9ZnHGgaZTG5vbj6yK8Paj+OVyHkcBxE8WuVfuB9ENWxf
Am+S8RVTm+rB3waLCcPfg4DMfXX9gDRyLxaH/txBVEue4qbz6QgiqwZSw3FXfVvu0STXjSwRsNJ3
Nve3UQn/y0RXMmCnrDvOeFkVojSAwEmdV0nmQwm5c5JLTuH8hbd6G9idR8p794FBsXFjQqjFjoSK
kyXs0mhOqiJMvSjyDI4oOzRYTCt7T/KAK3xj4tM9Alk3yeLlDg7wB7e5yIvUJpk/f+GDPIg2fJuJ
NVyFK6+umBPVYnj3k559KFAj0lUZSqe8jYB+guFFwr8oEjBWyv0qy8K2CDgDZfPZ85Bl52cXkB9w
c6g4t+f0YEGOCvLgKAbQuw8PDoHoxHRK5rz0HSujkX3X+zNzOLs4RgJeIxI6M4fzAQmqkS+UKKnK
QcdtCBdp67Qpe3XFtuwVw22meG7Jfe31/y8TMsSACQnUZ0Jc77xaioIr754kvYMNz8/uDR8LzVA1
NxY4O0Jb4dqW9u5uy/7NnsPO5gF87kdK3I95PIhMb+Jedrjbpvj7EfY+xdzd/t1s9nJ2z56yl+wZ
tLi4pOfzfZ7+/A8Us/OLFyG7fHkest4jfoZJRyY8sNjsMmTz1mQG3DPZfEz8qLpaJqWC3QrVD+Pu
9HjcdvB+yCm0OiTU7BKsvm2fklvwvHcKaNbcn0Hny+B7n2B+1hHgeUwwP3MEaOk8pb4KU1ihDHLY
qUF0X0XdpkQPooKuFB9TxWhybix7EO9fXbzTXV/pgo6m4xiyZchSFAvItocfXpiY4Nybie5Ke0fz
s0MCk15LpJfDNu8Z5WyHSmLtelm3GARia7IoqWteZr7NG4ojP3m6e7YMzNt9+PPovbO+AVEKoiww
b/dHw7WXpO3d/er8zZZDJb4iqBeIaPhofoE09pIVpiAeoz2uxDYukprgB9E2FNMrq1IguLXQuwM6
1RS+Zr9QyoRsipjDry0QeuDBEctMKExcqT4ixzfdYLD0XSQc8Pl+yAhBFKNlxNYgMZp+23t4BoYf
ZcO/u95UJKL0XXiKlWmkMZJ22azjuPOmTAQGyA87pXnxdiu07902pZk+qkazkxsGoU5ubhbDqbih
ibeb/iLP+qxISrFCN0W09vpl26K9FhwR0HOd2s42zud5vUlAfO23qHeeOTsYHBweIXrtHO5EaOft
KzuYdHwcoOVERsQkR7fdGUxMyoX1AT2RGzpaQqYJ2Ez+CENbJYadHoz+p2nA3NqNBIYTqnEJy39A
DJ/agB0O7Qc3vX/94cPCCw+OQ7bKG7W5ohjAfGKSGKMwGgwVwvn3vdqS05yqunJn1g01LH9mFELh
IRO56vdje4R7c47LIcWdkyOClriIG6V8Qu8aUVf0WuSmpkn7h2iOfZFsHcq4rQYoaX3YoGPtGfVM
0pazw3E/sge+dZ2bEo3VSPdxrHZSWtO2fK/9IW3g/ITV9GjM2vPO1AdpcO1bhO5iygnzfBxt5VCQ
NtHs+blnp60emw4+j16+mM69wX0tWn+gmh7OOFTTWkS0Om28RqgX5rIWJKtK92cyp6p1BOWSXcBa
y1twMEaFNAJVUMdFlWGPlnG3fhH73jbmeCCmSfwrN4weOqhvKueddINKOs4ae1zVWhRNEedmZSfT
TR3jREh4rilJm+k+r9risp+KTGTHhhks1aVPDUHFWmiFvAYHonPIFHP21SLQ+7XvEFxgHaZJb/Co
e4N2zX5jc/eEcZfcMZoWaG4haGbkz7oZah7abdqvg2ezYHRl3e5uLqFHCZ5kGUiPjU9Wr8E24o4C
CDw7JuFjm0TdTYCO1Z6Lm5O6cbE+Mi3Wbljc+zG0CNwVHN5KOByUCPVTaCepR1BdnxY0X3xaHJ2c
nL6+vzvZBkbQbswz1vIFVP0UHCXuBeSzq7ERYSkE2qseIKbvLes6bpAbFNqjqghovKSvCwqJU/ci
AABtMyJk8aOUPY/3eBQVfYLCpCTXUBKqzU5WwdOOMzth/vypKXYnMwNwko46xIDRb6z3heRIBpNp
To4zPu0u/0F9XwUHvA/6yjeXwAvKbg96r1FHF/2Qf2RiZi7BF+Z7zTDlHyUloztjqEbWGBasMsSN
Pj0NDBV8d9VPVzrJqZBi9rXahqbIDTqDQ7KdBfPYfOZ6iwGc9Oe1CP1iPr10nWVf77+5HmY1G7dU
2GvfBm0UO8NZmUbaez8Zw0B5MDodEh8MOIv+kABJRn15MWxKI36DIr8YBkfYgiuZcYkthhJVxWaY
NEN9l7hjJSWtFETsiirZD0w5YtEeHDHMIOpVG/ltFBxJijEP88E5NqYnDqBt3R62wHVeLZO8jTSD
4Tw/4qbSqianeK/zvP0I335VDX82CYcuucmGP4p+by2rB70JTThaodiyydZc/2o+7JvP2fSp+gta
AB0MPoX/mK37iB4ihNMGQhlx3J5/8hd9DM7dthR5tnXs96D+KtXuQg8SWttlyKxOWVPUynfhhDJP
C+nVPHjm/af0+kvDz5CHG0F/0Xj35/s/3mKPeP/2lr15e/vx3fW7N68/vrWbhflng3PDZldXMIdZ
9/Z7V/e/j4EvB+vHBK0sjssEoRlT5ntxTAtpHHsLtyrSdjr5L1BLAwQUAAAACAAAACtdcQXxhdoG
AACEEgAAEwAAAHZlcmlmeV9jYW5kaWRhdGUucHmNF9tu2zb0XV+hBSgk2bJquy3WZWawAUWGAUVT
FNtTFhCMRNtMdAtJ1YqL/vsOL5Io2c3mB0s6PPcbz7m4uPizzGhN4a+UPm1JKv0tK5mk/o5WBZX8
2Sdl5nMiWVWS3C8ACJjpnqaPIvG8mzJ/9uWe+kICHuGZn7N7ToCMCb8RNEv8DzRlBZDWvKorQYWf
VwfK/Vr/g+S6YqUUv3r0K1XS0pTWkmb9iWKkxQGsEazcDcrk1Y5wJvcFIKd5JRpOQaeLiwtvy6vC
x3jbSIBh7LOirrgEU8pKamrheR2M72rCBTU0mVXWHlrdYxCVkjytSklbGftfbv7+9AFff7y5+WLI
thw8p9h2hNcWYI+bMpVVlffHOW9wSsCsTosH0eEWRO47NCaeuDTgGsDg2u7kM3x2tJIV1PO8jG59
QTmjAoNjwjaWlBcCvVlHl54PPyLgVPqrDWo3aK1BRxS2i1X0Omznq8hAjug4O+pXHSBk3iW4LUed
UeHSIG8r7j/4rISIlDsaanlWWE80R5rP63A9e+iE9Nxn6Gj465RA65km0ZCmrgGi4fP1zPIAJloI
MJqFq8XxGBmGnEKcS8Ml1pTgj483f6yR45Be+zVQed5vfQzCgrSCHSn6VJU00n4EfHxfNWUmwnbs
vvZqqT8fkXke9iynfrtZDXa3M+tejbdAKxfxau0gvnYR5xYxr+I9cxVvjZFs6z9eoeVA3ls9f5wp
Y2+Xd0DZfazuvLNIKxdpeWcTB2qPbZ/DjEgSa2fHJK/3xNoOxc0lUnmWFBWUUFWyNDRa1RygAinC
28B8BHc2f1Rddyf6w54cWCb3uGwKCqVc8Q5nCrbY0JqqgpUKNEJ04BbThohBl1DtKKWhgxODApFu
Zg7QBtMSkjwPHeKDIVFZflBZPtUvcmlzWoaCytC4IIoQUhD71VdLrfgY4BDHU7XrQdn6aq2fSrf6
Vaa5ZEPNrWPdI8I6gpo4Uch43eryovrgj1ME6x24IOTeul6/v+BuOHUdDZ+nLu41GzupN1MrW/EM
jjUHEwEAaOdVfVDs9WBhY3smoSSDUssNIpv6HAMDiztJR1Zb1WJHKNGZoLSbOvubbOocRJ2qOyE8
J/t75wvHFPHUEE7R7TcyI6/qgYmJ/Cquo+/TpLKVB5cD/kryBgrz1oBotqM4hY4mbd/SvdsYhnUP
pzbwtNNhSNB7KjqyjpRZUvYfpLZxMYQexkD1U/cpKxs6OthVVYb+4hOoCc0T19HBzDweJmEyDus0
64yLTgVrV8b3HfmI5xl0awT5Cd2fP7QI4f2CRBAp6JCK9RP/MXZv6TXJBf0h2j2n5NE7I0qJUPTn
JZySAcl59CExuvvH/XUlvnmIw0DJJOlzmrNUmSfgos0ZTGXgNsqDmMUP0QkDnTpwv4bqJXYS85bd
OcPAAE9IrcZRjT+fdqTbB0sEHFnRFJqzQzxuaQZng9xLILTQ2AHa9jyMN/oCtLwGqLkQXQnLzXaz
Mm3rqkdbxeuom5+wgKmRotVyNlu/HQYprAVgM/ScG6oYzONtfCjHxXXSn4doHgY2h/LUNnPpQid1
x9gQ+iGM1rIdJwUAkprTFL1fjuBmfIcpMLRDMUiKXnfvrsgoycswmnVH26RXecDfJmeV7GfBwUwE
RRyGnfSEtjXw7n0bJbLCgEF3HPyqsyDkanCDRQE5Y3o0yJi4fSIuHlhPL2hnKNREEcxS/gb525lz
cojUGBYG4/XG3xKY/jIoEhXZaDwiu+kwNyOvHQR319jMwfoenImmCB1RQ75GWpXJXdBVCt+xEjlk
JzKVxgsYrQmsArNB6LiWFJerZewHH1UKdXug3MPKta/yzD9ALqkGAYH6Cltf4P3fnOvz7e27Eexc
FHuEjIk6J880w9Y+IXmflwZ0LuvsyShZ7ZQsmlyibwHMwZKlMMNeOrdxHGQMNIStFKt2CYdD14wD
vTWDHi3WnagXG1x2vWaU3mem10u3Q5nMCS6VRaYRxYFuPQZkutCEZaq03kIvVvsCNrGxnsE7uAtg
acJyT0qQNHHchJGgEKgMDNTeD6cT/0IvAvGb6Hs3+0NtqtU1yZqiFqFxo05yGDTW4+3MHNp9oyAw
cNoGptdvjrpFPPmd7xplwWcN7/YM/ZGQLMPEnodBCr2XwWBKg1g+1xSprfgF9MXC+Dbm9KmBiJo5
40UC43noqFuiEiRYJj+/e798u/oleJEMsgVuxROtAEUgS6IfikiEoxy0m5j2al4RKFiFk/SmJhDP
DJtiimJ9Zhc29epcUmpiUSCji7s2Kjm3jvPudPmMxTjdb2CSHDiT1Ah/Iezz4J8yUDs2qIBxSQqK
MUIBxirqGAeXtjOpFPD+BVBLAwQUAAAACAAAACtdnEkiMKIIAABBFQAACgAAAGNlcnRpZnkucHmV
WN1zozgSf89foZt9AByZ2EwyH3GUuqnNpfbhZie1s3cvPhdFQNiaYKCQcHD++v3pAxvHztUsD8Zq
dbda3T91t/jlHxetbC4eRXnByw2pt2pVle/P3r1799DwOmk4ScqMiDLjNcdPqYot2fBG5FuSkCZR
oiqTgkA8abZkXa3BQVLeKJGLNFE8PDv7r+bWA7CSVnJJqhJKHsxKniRSYYWkyUghHhtoCckdT8Ua
WjdJ0YJ9zRsOgbqp6krys7p65g3hZVpUsm24pOR5JdIVMbamKa8Vz+wSSa40Z5ekam9qUS1JWq2x
NSGrUoZn43Hd7zTLhOWCcMl5JknZrustUVW/OjaNF7g2nDzwpsGWNjxVVRNql52JdV01CqYsoRC2
5k21Jpnbjpt0u6OwJE2KtCoV7xQlf3z7z+938f2/v337w4rlDeyGNbIXvHcEkkhy73jaEotXxY6n
aNo4TdIV701ZJXIFx/bDH9izFa0TpSd6wQcM7QSE06d4yRFLhZi6eUOlRK6rJx4rLtXZ2VnGc+3O
WGI7hd/Rkn2cBNdnBM8L87vxNLjwu/NpMHuJ2MvoZabYy0yye38SGJ68asgPQAuxKZfcL52ofuQ5
i0bq4t6PRj+0AjViL5GZbbhqm5JIKs81h69ZSrCM/On4JQoCmPXPnRP836uSBzs7BVzdAFR+55bq
YEwXzBIpgVjS3U5mT2xiZgCpgpPuZnrdjVg0exqz6ZB+G113F5p+7uhFRVeCDZwRzIqIttGABEuD
YLgHv6jOn0ZgWwm82yggIidPt2xCeAGk2enWTRdR4Dzu4AR/A7epkEAE+3DldvQs1OoAWH6g4ZKq
bu9cDEItynbyB1NN1ZaZKJdsiMiewZl+19sQ4oBwnK2qCS72ROSJao2UoMlvBcQc40FIaH4YlByv
3PoLbukYm167xe/9aUD1z9/f8afJjrZl/t/YRTAa7djzU+z5IftuGV5L1vN4Uz6+uvIADb137G87
xnQwa+vaDs/10DqhoF3LXoF2VhS0eEU1qqCioO2rGaPVanMAL9oblo+g2vfyBDhGynW5NKsryHlU
R2Eo0Ra3WqLdSxitpyRccIxGargcXm218LNEJS7CKRCmJDMpxdJndQNtvvcvk6t3uefhy/fv1x61
/DQvWrlifzYtd/kDLtPScy/3FjjFu2FihkW9GpD0SJPXzBJQz9TKWww3O7nJb6am3CW37q2lbgE1
isNp6qAUpa5XKffXFBZb6vqWRUNFO7OAhiJ+FplaxRIA5TCAMavO8G/YXGPdZMJOZ0IrZyuKt1jM
pHTmYlEkXW/R56qCl/4GyvRbSmeG0LRbC/GmqpST1X/3koMt6In9LiY3TBNuhio3c01aMJflnviW
SWlpc+8ZXUH1PDAK0/MJmPU7wntidOjR1FHfu/eldfwahRLWsPliVxAEldoR3B0xrk3ZH+RiNagf
vcijFpBzD91DCUzBoL2AflqIjB7hQxeCGbScs1f5pwWOYfxoMwdniiyfwf87Nc7QMKl1F+RDwTgB
q1jsiofjmI+ni9vJ/rg01bNHBe1z9oDL5QgH/K+iFOt27ZqVJcTUaoxfghanBu7JrbdXgjg7RQGN
PgVHB6NIaJscJgMcseKStpeH1Hv/Mghmy6RmRTI2YB+1l0MkYwpF0dkIUeixfLqeXg6NAudJW1Jm
ze2wmBxgBpAYH4wnC3owvlwE9CCM9tHR1wXhgDda3Lii+WqR968WiY4XCS50HAcrvYVC59uUtumh
F9NhrR/4ti/3v5DfkgL9okkHxKQDdL0V4ciMW0ftxXS/ZlnQFzci1e32I0dSRUeLZlqrM5krtgBg
RXrur9FojXYBXCOE4xx9wzCOQ5l9QO9FB4iaOdcrr2AoSr8zEql3TXXY+8b+BBiHmk8CwNWFDHvx
bRZjNp/riPdHtidNF7RP/vGgSe+nEb4jSABeMQ5K7GyLbW2F8/y3zssJXOnI2UP3f/Q4iJ8Qx+FV
IinEi7lkxKkRSwM6dM4Jha99F/QtnruQmDJCc5pQE1qrzmVD15bb+wkwU9aGfFwwyr5GzL6ysg5f
eFNJ3y9pGQQ/mXTfyLFf54IOcuU5y4sqUf5hrkXf1JNztyBQLilKnLam0OBfhlws/a8Bmm+Q4A2d
LTRX2PCkCGYbBu75NX1aGMJsc8E282Fps6f3jk0no9E0mm1QVrWb733Tyvrd6C6gd8GgzG4WfV+Z
MOYlraq864RZGX0yzOpzt94YXdunwOk4CelWMl3rheJHN+LBVXhGmrZ0hG1ob5W2N/KOEWW7Fxv2
nAECzIWfmV96VN6pbRrYZkPdKZPyWO1RR8K86UWE0yxkWul01J8CgwSRupgO3WFACs/FcZmseRzD
gTEwDIDHnsUMykl/AQ6/NMtWp44HPWp8lEtEOMvixNF9T5R1iz5SbWvO9CX0BMvuhu5RewNmgHfV
4B6KLOO5hvVYqGrVz6jWx1LnszxpC8W8STi9uvx89fnz5G3N1kkDoSj89GEyff/x40fv5CK2/Rwu
8vHq0+Qyev/2GrZBtcYDkzvZ6aR/Tq5kLuhjfUE/7awZmCWDoImPFpW+tUGjienvAyFinklfT4Um
OjruWWzvVrsbmZl2cdnnCldvzKT1PxpOUlaK6Jvfjs2sdZDlrDpzeTB/rYvtf3sMzF+bAvfL7dcJ
nxscQGul2UWG1CitcmG+XLEoOPf+V3r9IdYJlA1vJwdb23/puB7cfrUMLg8Cl0kVr6usLZIm3rN6
C7YfDJQOOr07I0ucrP2gQjR7f9v5iUWOimxvtHXF9Wn/CFPd8lx0vhduBl/lQu0vL3jDg9aekz5U
0M8VO4j2ITJsV7aH0rDn/XWfG8n3376Mo6sP2L77YBXKVQKCb5ewCHzcYvuAYLjiXSaW2sXH/YZT
bpzpvi0WgN4uIaOrwLBUfU+jVg2Xq6rAbU+R/pwO74zU07eYvKfmr9z/F1BLAwQUAAAACAAAACtd
qBvvpMgEAACuCwAAEQAAAGNoZWNrX2dlb21ldHJ5LnB5fVZLc9s2EL7rV7DpeAjIkC25PnQkwdNe
PFMf8pgcfOCwHIiELCQkSANgJCWT/97Fg6RkKtVBJPf57WLxAb//dttqdbsR8pbLb1FzNLta/jF5
9+7dP7LgDYc/aSJ+YLmJXnhdcaOOUb7j+VcdbWsVKWZELVkZvShRzJq6FPkxevr84X20FSXXNxBo
slV1FW0VhABTHYmqqZWJHoMgYjp6DDatzE1dl71NqdosZ5BuEgRfdC29bcPMrhSbzvIjfHZG+qgn
k0nBtx4oKphheDmJ4KcNM1xTK0li/xGnq72QRb3XNEmdka0LIMhg7T3tj5ENyUlBXqlO4uAUp4Na
aw7phRYSPGXO0SsR0uCIySJ6fZj3hgKQrefLYJ9TWlA6W/TqZ5qgR8TIKyaPaAMPTN7Xkge3iJea
R2CQe4PCGgwgLPi9Bf884A45QawjWZvIRuuyz9d0n8zT9T5ZpGs6gAjl3bDGDgF6xk6zUVDXDjo4
XwnJX1tWCiPsZ984QVzruGwrDsPBkW8iHtA80R8/V4prUbTgqblB+Az9xrU+ibtccXpeiaKbJFaw
bKECZe3B9oWje99r5aq00pClT3fDigIpfBYvhIGges+aOLWOaE4WPpbNxbclz82J5pI/NBJs850o
izhdl1x2lZ8Zt/QRTW2qnJU8TvHKdIKd2BorCOHak3k57fz1yRp1HWuGDtzh82a5paRhMZMBYJo0
//YVpyOXblrcpOS1NEK2/JLVaXuWe4oWMztIxD7mKR55NCXLeUGRuW6n1oL4lwXU3TAFVNNDFWmi
ru7GyEJ7vPXpQP+/pR1x6rPb1/BmRz7oATQSsL8b4nXEy8clPCXq+n7apCHa2T64pndno/yFuK34
dCMMrzTCoy2Jvlwv8NWfzmj5qwr81nxKgnFq+wZYITgZhGOg58D81ChuWiWjk+kk3WCRU/vAn7qq
v/IMzIwjUQIngtnR+1DHmO+cfuA89/kQUl+kXqf5q6d5ZBfSV2LTb1qYVaRJhS9ysIuR6PQSF0Nv
K0oXo4bnD3Tusc3yh82MLUNHEFqQc0ol5x3t7ebkDTcPdsAybWngGFl9ovfTKapmizfU5no9HC0O
/K94zlOFU474opMG0hgxwoGUNdkJm8k3sfMI259YaOOBO1KEEGimfYZ+c+Pp4bpP61gDX326yAhv
PZcOCl3MdgJ4oaxHTr5t3UnTQwWKh412dBQBIexjJzB+e9zaWf5xCFVnJAu8DxF/Ykqt1n+dLpNb
StM2JT9TsrLZsQ031PfMj6qqayjCjz5eVXXRlq22y+skKw17BmLQH9+n36+C1oH5PpByEOOfLkvD
hLIHKFPKXjqGo/NAYLitV4djORmfMYd1l8ROMQjYenN6bttAR7tDLgcKi3Sg9HiZ1x26t4eMZarj
7ID7Au0E+8J/yVqAKifoQI4kbFkcKn7DRYXIjScO6umja3F4kkK8CKPdSnb1YFKrgiteZL6Z7p94
RFnoq39gYDJAn2WSVTzLKI2zrGJCZlm87Btvda6kI1wQ1Mu3ZLE82Yt53UrI7++S9gJ6U9as0Mhe
OZF1xTeKswJo8gBXmZMBbRQwobMg8ce/P39ehl1P+sOchJu1XQkAVcGBc+sqA0kFB5CAbHoZEw8B
T/4DUEsBAhQDFAAAAAgAAAArXVK9U9hDAQAAeAIAAA0AAAAAAAAAAAAAAIABAAAAAG1hbmlmZXN0
Lmpzb25QSwECFAMUAAAACAAAACtd3Fobg7wBAAAUBwAADgAAAAAAAAAAAAAAgAFuAQAAZGF0YS9x
MjE1Lmpzb25QSwECFAMUAAAACAAAACtddj8e34kPAAAZHAEADgAAAAAAAAAAAAAAgAFWAwAAZGF0
YS9xNDM3Lmpzb25QSwECFAMUAAAACAAAACtd4aW9R1wCAADFHAAAEAAAAAAAAAAAAAAAgAELEwAA
ZGF0YS9iaW5hcnkuanNvblBLAQIUAxQAAAAIAAAAK109JQUlrAEAAIQKAAARAAAAAAAAAAAAAACA
AZUVAABkYXRhL2xpYnJhcnkuanNvblBLAQIUAxQAAAAIAAAAK11tQ+nkSQkAAJwZAAAJAAAAAAAA
AAAAAACAAXAXAAB2ZXJpZnkucHlQSwECFAMUAAAACAAAACtdcQXxhdoGAACEEgAAEwAAAAAAAAAA
AAAAgAHgIAAAdmVyaWZ5X2NhbmRpZGF0ZS5weVBLAQIUAxQAAAAIAAAAK12cSSIwoggAAEEVAAAK
AAAAAAAAAAAAAACAAesnAABjZXJ0aWZ5LnB5UEsBAhQDFAAAAAgAAAArXagb76TIBAAArgsAABEA
AAAAAAAAAAAAAIABtTAAAGNoZWNrX2dlb21ldHJ5LnB5UEsFBgAAAAAJAAkAHwIAAKw1AAAAAA==
"""

def main():
    if not __debug__:
        raise SystemExit('Run without -O or -OO: verification requires assertions.')
    parser = argparse.ArgumentParser(description=__doc__)
    parser.add_argument('--extract', type=pathlib.Path)
    args = parser.parse_args()
    payload = base64.b64decode(PAYLOAD)
    if hashlib.sha256(payload).hexdigest() != ARCHIVE_SHA256:
        raise SystemExit('Certificate archive checksum mismatch.')
    def verify(destination):
        destination.mkdir(parents=True, exist_ok=True)
        with zipfile.ZipFile(io.BytesIO(payload)) as archive:
            for entry in archive.infolist():
                relative = pathlib.PurePosixPath(entry.filename)
                if relative.is_absolute() or '..' in relative.parts:
                    raise ValueError('Unsafe archive member')
                target = destination.joinpath(*relative.parts)
                if target.exists():
                    raise FileExistsError('Use a fresh extraction directory: '+str(target))
                target.parent.mkdir(parents=True, exist_ok=True)
                target.write_bytes(archive.read(entry))
        subprocess.run([sys.executable, str(destination/'verify.py')],
                       cwd=destination, check=True)
    if args.extract:
        verify(args.extract.resolve())
    else:
        with tempfile.TemporaryDirectory(prefix='sarkozy-simple-') as directory:
            verify(pathlib.Path(directory))
if __name__ == '__main__':
    main()
\end{filecontents*}

\title{A simpler construction of square-difference-free sets\\beyond exponent $0.758$}
\author{Eric Naslund}
\date{September 11, 2026}
\begin{document}
\maketitle
\begin{abstract}
We construct subsets of $\{1,\ldots,N\}$ with no nonzero square difference
and cardinality at least $N^{0.758001-o(1)}$. The construction improves on
Krachun's exponent $0.752796455874514\ldots$ and is a simpler companion to
our construction of exponent $0.75806746$. Its purpose is to make the
mechanism accessible: allow modular square differences, order them with
intervals, and combine the resulting local objects by the Chinese remainder
theorem. A short stopping-tree argument turns interval moments into an
exponent. The concrete ingredients are six small prime chains, two odd
components with only 17 and 729 low points, and an explicit seven-state
binary recursion. All general implications are proved here. The remaining
finite checks use exact rational arithmetic; their complete data and a
standalone verifier accompany this paper in its source and PDF attachment.
\end{abstract}

\section{Introduction}
For $N\ge1$, let
\[
 D(N)=\max\bigl\{|A|:A\subseteq\{1,\ldots,N\},\quad
 (A-A)\cap\{1,4,9,\ldots\}=\varnothing\bigr\}.
\]
The Furstenberg--S\'ark\"ozy theorem says that $D(N)=o(N)$
\cite{Furstenberg1977,Sarkozy1978}. We study the complementary question:
how large an exponent $\alpha$ can a construction achieve in
$D(N)\ge N^{\alpha-o(1)}$? This notation means that for every
$\varepsilon>0$ the inequality $D(N)\ge N^{\alpha-\varepsilon}$ holds
for every sufficiently large integer $N$.

\begin{theorem}\label{thm:main}
With $\al=758001/1000000=0.758001$, we have
\[
 D(N)\ge N^{\al-o(1)}.
\]
\end{theorem}

\paragraph{Previous constructions.}
Ruzsa obtained the exponent
\[
 \frac12+\frac{\log7}{2\log65}=0.733077\ldots
\]
by an alternating-digit construction \cite{Ruzsa1984}.
Beigel and Gasarch, and independently Lewko, obtained
\[
 \frac12+\frac{\log12}{2\log205}=0.7334117970\ldots
\]
using a 12-element square-difference-free set modulo 205
\cite{BeigelGasarch2008,Lewko2015}.
The often-mentioned three-quarter barrier concerns this particular
modular-avoidance strategy and conjectured limitations on its inputs;
it is not an upper bound on $D(N)$.

Krachun's recent construction crosses that barrier
\cite{Krachun2026}. Instead of requiring local residue sets to avoid all
square differences, it permits such differences when they consistently
increase a rank (in our sign convention). Ranks from coprime components add, while the numbers of
digits multiply. His explicit prime-chain inputs give
\begin{equation}\label{eq:krachun}
 \alpha_{\rm K}=
 \frac{10+\sum_{i=1}^{10}\log p_i/\log t_i}
      {1+2\sum_{i=1}^{10}\log p_i/\log t_i}
 =0.752796455874514\ldots,
\end{equation}
where
\[
 (p_i,t_i)=(3,2),(7,3),(11,4),(19,5),(23,5),
 (31,7),(43,7),(59,9),(71,9),(103,11).
\]
Our theorem raises this exponent by approximately $0.005204544125486$.
The ranked realization, repetition, and CRT mechanism are due to Krachun;
we prove the needed statements below, with the reversed rank convention
\cite[Lemmas 4--5]{Krachun2026}.

For perspective, the upper-bound problem has developed through the work of
Pintz, Steiger and Szemer\'edi, Bloom and Maynard, and Green and Sawhney
\cite{PintzSteigerSzemeredi1988,BloomMaynard2022,GreenSawhney2025}.
Green and Sawhney prove $D(N)\ll N\exp(-c\sqrt{\log N})$ for an absolute
$c>0$ \cite[Theorem 1.1]{GreenSawhney2025}. There remains a large gap
between these upper bounds and the construction problem considered here.

\paragraph{The extra flexibility.}
Replace a discrete rank by an interval. If $y-x$ is a square modulo the
local modulus, require the interval of $x$ to lie to the left of that of
$y$. Intervals for incomparable digits may overlap, and their lengths
need not agree. Repeating digits multiplies interval lengths; this makes
\emph{moments of the lengths} the natural measure of a local construction.
For a target exponent $\alpha$, a component in a square base $b$ earns a
contribution $f$ when its widths satisfy
\[
 \sum_x w_x^f\ge b^\alpha.
\]
Section~\ref{sec:criterion} proves that coprime components achieve the target
as soon as their contributions add to more than $\alpha$. The proof uses
a finite tree whose mass is conserved at every branching. No limiting
frequency argument or optimization theorem is required.

The two additional local ideas are to group odd primes and to use the
prime 2. For odd primes, restricted digits are interleaved with free digits;
the first differing position is inspected separately at each prime.
For 2, a state records where the even and odd digit intervals may lie.
The fact that odd squares are 1 modulo 8 then leads to a small recursion.
The moment method accommodates both kinds of inputs in the same argument.

\paragraph{Why this companion is simpler.}
A separate manuscript attains $0.75806746$ with larger finite witnesses
\cite{NaslundCompanion2026}. Here the target is deliberately lower:
the two odd supports have 17 and 729 points, instead of 4,913 and 19,683,
and the binary recursion has seven states and 26 branches, instead of
25 states and 94 branches. Its windows, transitions, and positive comparison
vector fit in explicit tables. The final budget is the exact addition
\begin{equation}\label{eq:budgetintro}
 0.51011+0.0268+0.0673+0.154=0.75821>0.758001.
\end{equation}
This paper proves its own result without using the stronger companion.

The theorem remains computer-assisted: the 729-point component carries a
finite table of widths. We specify how those data define the construction,
prove why the verification conditions suffice, and supply the complete
checks. Optimization helped find the inputs; no optimizer or floating-point
comparison is needed to accept them. All logarithms are natural, and all
terminating decimals specifying parameters denote exact rationals.

\section{From local order to square-difference-free sets}\label{sec:ranked}
The local goal is an order compatible with square differences, rather than
an immediate square-difference-free set. The final lift will put a higher
rank into an earlier integer block, turning every potential positive square
difference into a negative number.

For $b\ge1$, let
$\QR(b)=\{z^2\bmod b:z\in\Z\}$, including zero. Residues are represented
canonically in $\{0,\ldots,b-1\}$. For distinct $x,y$ write $x\to_b y$ if
$y-x\in\QR(b)$. For composite $b$, a nonzero square residue may vanish
in some prime coordinates; invertibility is never imposed.

\begin{definition}
A ranked alphabet modulo $b$ is a finite set $C$ of canonical residues
with an integer function $r:C\to\{0,\ldots,h-1\}$, for some $h\ge1$,
such that $x\to_b y$ implies $r(x)<r(y)$.
An ordered interval alphabet assigns to every $x\in C$ an interval
$[a_x,a_x+w_x]\subseteq[0,1]$, with $w_x>0$, such that
$x\to_b y$ implies $a_x+w_x\le a_y$.
Its moment of order $f\ge0$ is $Z_f(C,w)=\sum_{x\in C}w_x^f$.
\end{definition}

Intervals may overlap when the corresponding digits are incomparable. The
graph of an ordered interval alphabet is acyclic, since a directed cycle
would imply that the sum of its positive widths is at most zero.

The small abstract example in Figure~\ref{fig:intervals} illustrates the
flexibility. If the only required arrows are $a\to b\to c$ and $d\to c$,
the interval for $d$ can occupy two thirds of the unit interval. This is
compatible with the shorter intervals for $a$ and $b$. The same idea applies
to the directed modular-square graphs used in our construction.

\begin{figure}[ht]
\centering
\setlength{\unitlength}{1mm}
\begin{picture}(112,37)
\put(5,21){\makebox(0,0){$a$}}\put(20,21){\makebox(0,0){$b$}}
\put(35,21){\makebox(0,0){$c$}}\put(20,6){\makebox(0,0){$d$}}
\put(8,21){\vector(1,0){8}}\put(23,21){\vector(1,0){8}}
\put(23,9){\vector(1,1){9}}
\put(60,28){\line(1,0){45}}
\put(60,27){\line(0,1){2}}\put(75,27){\line(0,1){2}}
\put(90,27){\line(0,1){2}}\put(105,27){\line(0,1){2}}
\put(60,32){\makebox(0,0){$0$}}\put(75,32){\makebox(0,0){$1/3$}}
\put(90,32){\makebox(0,0){$2/3$}}\put(105,32){\makebox(0,0){$1$}}
\put(55,22){\makebox(0,0){$a$}}\put(55,16){\makebox(0,0){$b$}}
\put(55,10){\makebox(0,0){$c$}}\put(55,4){\makebox(0,0){$d$}}
{\color{blue!65!black}\linethickness{.7mm}
\put(60,22){\line(1,0){15}}\put(75,16){\line(1,0){15}}
\put(90,10){\line(1,0){15}}\put(60,4){\line(1,0){30}}}
\end{picture}
\caption{Intervals may overlap when no order is required. The moment here is
$3(1/3)^f+(2/3)^f$. This is an illustrative directed graph, not a separate
arithmetic input.}\label{fig:intervals}
\end{figure}

\begin{lemma}[Path budgets]\label{lem:paths}
Let a finite directed acyclic graph have positive vertex weights $w_x$.
There are intervals in $[0,1]$ of those widths ordered along every edge
if and only if the sum of the weights on every directed path is at most one.
\end{lemma}
\begin{proof}
Ordered intervals along a path have disjoint interiors, so their widths
sum to at most one. Conversely, let $d_x$ be the maximum weight of a path
ending at $x$, including $x$. Set $a_x=d_x-w_x$. Then $a_x\ge0$ and
$a_x+w_x=d_x\le1$. If $x\to y$, appending $y$ to a maximizing path ending
at $x$ gives $d_y\ge d_x+w_y$, whence $a_y\ge a_x+w_x$.
\end{proof}
This is the weighted-chain viewpoint underlying the chain polytope of a
finite poset \cite{Stanley1986}; the elementary proof above supplies the
part used here.

\begin{lemma}[Reverse-rank realization]\label{lem:reverse}
A ranked alphabet of size $q$ modulo $M$, with ranks in
$\{0,\ldots,h-1\}$, gives a square-difference-free set of size $q$ in
$\{1,\ldots,Mh\}$.
\end{lemma}
\begin{proof}
For the canonical residue $x$, choose
\[
 T(x)=1+x+M(h-1-r(x)).
\]
These numbers lie in the stated interval and are distinct modulo $M$.
If $T(y)-T(x)=z^2>0$, then $x\ne y$ and reduction modulo $M$ implies
$r(y)-r(x)\ge1$. Consequently
\[
 T(y)-T(x)=y-x-M(r(y)-r(x))\le(M-1)-M<0,
\]
a contradiction.
\end{proof}

\begin{lemma}[CRT gluing]\label{lem:crt}
Suppose $M_1,\ldots,M_n$ are pairwise coprime and $C_i$ is a ranked
alphabet modulo $M_i$, with $q_i$ digits and ranks below $h_i$.
Their CRT product is a ranked alphabet modulo $M=\prod_iM_i$, with
$Q=\prod_iq_i$ digits and ranks below
\[
 h=1+\sum_{i=1}^n(h_i-1).
\]
\end{lemma}
\begin{proof}
The Chinese remainder map identifies residues modulo $M$ with tuples of
residues modulo the $M_i$. Explicitly, writing $P_i=M/M_i$, choose
$u_i$ with $P_iu_i\equiv1\pmod{M_i}$ and take the canonical residue of
$\sum_i x_iP_iu_i$. This proves existence; divisibility by every $M_i$
proves uniqueness. Assign rank $r(x)=\sum_ir_i(x_i)$. A square difference
modulo $M$ is a square in each coordinate. Equal coordinates leave the rank
unchanged and unequal ones increase it. At least one coordinate is unequal,
so the sum strictly increases. The size and rank bounds follow.
\end{proof}
Lemmas~\ref{lem:reverse} and \ref{lem:crt} express the ranked realization
and gluing mechanism of \cite[Lemmas 4--5]{Krachun2026} in our convention.

\begin{lemma}[Square-base words]\label{lem:words}
Let $b=s^2>1$. If digits $C$ form a ranked alphabet with rank height $h$,
then their length-$k$ words form a ranked alphabet modulo $b^k$ with
$\card C^k$ digits and rank height $h^k$. If the digits instead carry
ordered intervals, the word intervals obtained by affine composition are
ordered along modular square differences as well.
\end{lemma}
\begin{proof}
Encode a word by $X=\sum_{j=0}^{k-1}x_jb^j$. Uniqueness of the base-$b$
expansion gives injectivity. Suppose $Y-X\equiv z^2\pmod{b^k}$ and $X\ne Y$.
Let $j$ be the first position at which $x_j\ne y_j$, reading from the
least significant digit. Since $s^{2j}=b^j$ divides $z^2$, we have
$s^j\mid z$: this follows by comparing the valuations of every prime
dividing $s$. Dividing the congruence by $b^j$ and reducing modulo $b$
gives $y_j-x_j\equiv(z/s^j)^2\pmod b$, a nonzero square edge.

In the ranked case assign the word rank
$R(X)=\sum_{j=0}^{k-1}r(x_j)h^{k-1-j}$. The first differing digit increases
its rank, and its place value dominates every possible decrease in the
remaining places. Thus $R(Y)>R(X)$. Its range is $0\le R<h^k$.

For intervals, put $F_x(t)=a_x+w_xt$ and assign
$F_{x_0}\circ\cdots\circ F_{x_{k-1}}([0,1])$. At the first differing
position, the two digit intervals are ordered. The remaining compositions
lie inside those digit intervals, and the common preceding maps are
increasing. This preserves the required order. The width is
$W_X=\prod_{j=0}^{k-1}w_{x_j}$.
\end{proof}

\begin{lemma}[Rounding intervals to ranks]\label{lem:round}
If every interval in an ordered alphabet has width at least $\gamma>0$,
then $r(x)=\lfloor a_x/\gamma\rfloor$ is a valid rank with
$0\le r(x)<\lceil1/\gamma\rceil$.
\end{lemma}
\begin{proof}
On an edge $a_y\ge a_x+\gamma$, so $r(y)\ge r(x)+1$.
Moreover $0\le a_x\le1-\gamma$, which gives the stated upper bound.
\end{proof}

\section{A finite interval-moment criterion}\label{sec:criterion}
\begin{theorem}\label{thm:criterion}
Let $n\ge1$. For each $i\in\{1,\ldots,n\}$ let $C_i$ be an ordered
interval alphabet in a square base $b_i>1$, with the bases pairwise coprime.
Suppose there are common constants $0<\sigma\le\rho<1$ such that
$\sigma\le w_{i,x}\le\rho$ for all digits. If $\alpha\ge0$, $f_i\ge0$,
and
\[
 Z_{f_i}(C_i,w_i)\ge b_i^\alpha\quad(1\le i\le n),
 \qquad F:=\sum_{i=1}^nf_i>\alpha,
\]
then $D(N)\ge N^{\alpha-o(1)}$.
\end{theorem}
The reason for the strict surplus is simple. At scale $\delta$, the
number of available words is controlled by $\delta^{-\sum f_i}$,
whereas the common rank height costs only a factor of order
$\delta^{-1}$. The strict inequality $\sum f_i>\alpha$ pays for this
height and for selecting a common length within each component.

The common width bounds are automatic for a fixed finite collection of
nonempty alphabets with every width strictly between zero and one. They need
not be uniform over a sequence of proposed constructions.

\begin{lemma}[Stopping at a common width]\label{lem:stopping}
For one alphabet in base $b$, suppose $\sigma\le w_x\le\rho<1$ and
$Z=\sum_xw_x^f\ge q=b^\alpha$, with $f\ge0$.
For every integer $K\ge1$, put $\delta=\rho^K$. There are a length
$1\le k\le K$ and a set $D$ of length-$k$ words such that
\[
 \sigma\delta<W_X\le\delta\quad(X\in D),\qquad
 q^k\le K|D|\delta^f.
\]
\end{lemma}
\begin{proof}
Starting at the empty word, extend a word by every digit until its width
first becomes at most $\delta$. Every branch stops by depth $K$, since a
length-$K$ word has width at most $\rho^K$. A stopped word has a parent
of width greater than $\delta$, so its width is greater than
$\sigma\delta$.

Give the edge labelled $x$ the probability $w_x^f/Z$. These probabilities
sum to one. The mass of a word $X$ of length $k$ is therefore $W_X^f/Z^k$.
Replacing a node by all its children preserves its mass, so the terminal
masses in this finite tree sum to one. These are masses on words;
overlap between their geometric intervals causes no difficulty.
Let $T_k$ be the sum of $W_X^f$ over the words stopping at depth $k$.
Because $Z\ge q$,
\[
 1=\sum_{k=1}^K\frac{T_k}{Z^k}
 \le\sum_{k=1}^K\frac{T_k}{q^k}.
\]
At least one summand on the right is at least $1/K$. Retain the words
stopping at that length. They have $W_X\le\delta$, and hence
$q^k/K\le T_k\le |D|\delta^f$, as required.
\end{proof}

\begin{proof}[Proof of Theorem~\ref{thm:criterion}]
Apply Lemma~\ref{lem:stopping} separately in each component with the same
$K$ and $\delta=\rho^K$. Obtain lengths $k_i\in[1,K]$ and word sets $D_i$.
Their moduli $b_i^{k_i}$ are coprime perfect squares, and
Lemmas~\ref{lem:words} and \ref{lem:round} give ranks below
$H=\lceil1/(\sigma\delta)\rceil$ in every component.
Put $M=\prod_ib_i^{k_i}$ and $Q=\prod_i\card D_i$.
Lemma~\ref{lem:crt} gives a ranked alphabet of size $Q$ modulo $M$ with
height $h=1+n(H-1)$. Multiplying the selected-moment estimates yields
\begin{equation}\label{eq:productmoment}
 M^\alpha\le Q\delta^F K^n.
\end{equation}
Since $\delta\le1$, the fixed constant $c=(n+1)(1+\sigma^{-1})$ satisfies
$h\le c/\delta$. Consequently
\[
 (Mh)^\alpha\le Qc^\alpha K^n\delta^{F-\alpha}.
\]
Because $F-\alpha>0$ and $0<\rho<1$, the last factor
$c^\alpha K^n\rho^{K(F-\alpha)}$ tends to zero.
For example, the ratio of consecutive terms tends to
$\rho^{F-\alpha}<1$, which proves this without any asymptotic counting
formula. Choose and fix $K$ so that $(Mh)^\alpha\le Q$.

Write $L=Mh>1$. Repeating this one ranked alphabet $t$ times and applying
Lemmas~\ref{lem:words} and \ref{lem:reverse} gives $Q^t$ square-difference-free
integers in $[1,L^t]$. For any $N\ge L$, let $t$ be the integer with
$L^t\le N<L^{t+1}$. Then
\[
 N^\alpha\le L^{(t+1)\alpha}\le Q^{t+1},\qquad
 D(N)\ge Q^t\ge Q^{-1}N^\alpha.
\]
For every $\varepsilon>0$, $N^\varepsilon\ge Q$ eventually, giving
$D(N)\ge N^{\alpha-\varepsilon}$ for every sufficiently large $N$.
All alphabets, $K$, and hence $L,Q$ have been fixed before $N$ varies.
\end{proof}

\section{Odd-prime digit constructions}\label{sec:odd}
\subsection{Prime chains}
\begin{lemma}\label{lem:chain}
Let $p\equiv3\pmod4$ be prime and let $(s_0,\ldots,s_{t-1})$ be distinct
residues such that $s_j-s_i$ is a nonzero square modulo $p$ whenever $i<j$.
Then the $pt$ digits $s_j+pu$, $0\le u<p$, form an ordered interval
alphabet modulo $p^2$, with interval $[j/t,(j+1)/t]$ on each digit
$s_j+pu$. Its moment is $pt^{1-f}$.
\end{lemma}
\begin{proof}
The element $-1$ is not a square modulo $p$: otherwise an element of order
four would exist in the multiplicative group of order $p-1$, contradicting
$4\nmid p-1$. A square difference between different low digits therefore
points from the earlier to the later chain entry. If the low digits agree,
a nonzero difference has valuation one, and cannot be a square modulo $p^2$.
The intervals are ordered in the first case and no edge is possible in the
second. There are $p$ free copies of each interval, giving the moment.
\end{proof}
This is the local Paley-chain mechanism of \cite{Krachun2026}; the
alternating free-digit principle goes back to \cite{Ruzsa1984}.

\subsection{Separate first differences at different primes}
Fix distinct odd primes $p_1,\ldots,p_s$ and positive integers
$e_1,\ldots,e_s$. A low point $a$ consists of words
$(a_{p,0},\ldots,a_{p,e_p-1})$ with $0\le a_{p,j}<p$.
For distinct low points $a,b$, impose an edge $a\rightsquigarrow b$ if,
for every prime $p$, either the coordinate words are identical or their
first differing position $j$ satisfies
\begin{equation}\label{eq:lowedge}
 b_{p,j}-a_{p,j}\in\QR(p)\setminus\{0\}.
\end{equation}
The first differing position is allowed to depend on $p$. No comparison is
required from a prime whose entire low word is unchanged.

\begin{lemma}[Prime-coordinate free-digit lift]\label{lem:odd}
Suppose a finite set $S$ of distinct low points has positive intervals in
$[0,1]$ ordered along every edge $\rightsquigarrow$.
Put $F=\prod_pp^{e_p}$. For each $a\in S$, choose arbitrary
$u_{p,j}\in\{0,\ldots,p-1\}$ and form
\begin{equation}\label{eq:oddlift}
 x_p=\sum_{j=0}^{e_p-1}
       \bigl(a_{p,j}p^{2j}+u_{p,j}p^{2j+1}\bigr)\pmod{p^{2e_p}}.
\end{equation}
Combining the $x_p$ by CRT produces an ordered interval alphabet in base
$F^2$, with $F\card S$ digits. Every free copy uses its low point's interval,
and the moment is $F\sum_{a\in S}w_a^f$.
\end{lemma}
\begin{proof}
Unique prime-adic digit expansions and CRT show that all these digits are
distinct and count their number. Suppose different lifted digits $x,y$
have a square difference modulo $F^2$. Consider a prime coordinate in which
they differ. If $k<2e_p$ is the first differing actual digit position, then
$y_p-x_p$ is divisible by $p^k$ but not by $p^{k+1}$: all lower digits
agree and the coefficient at $p^k$ is nonzero modulo $p$.
The congruence to a square forces $k$ to be even. Indeed, if
$z^2\equiv y_p-x_p\pmod{p^{2e_p}}$, the valuation of $z^2$ below
$2e_p$ must equal $k$, so $k=2j$; dividing shows that the leading nonzero
difference is a square modulo $p$.

If the low coordinate words differ, the first actual difference must thus
occur at their first differing low digit, and \eqref{eq:lowedge} holds.
If two distinct lifted digits have the same low point in every coordinate,
some first actual difference would instead lie in a free, odd position,
a contradiction. Thus their low points are distinct, and each coordinate
is either identical as a low word or obeys \eqref{eq:lowedge}. The assumed
low interval order therefore gives the required lifted order. Counting
free copies gives the moment formula.
\end{proof}

The first differing position must be tested separately at each prime.
For example, the 19-coordinate may first differ in position 0 while the
23-coordinate first differs in position 1. If one coordinate word is
unchanged, its zero difference is allowed. The lemma therefore does not
require the difference to be invertible modulo the composite base.
Only the explicit low points and their interval order enter the proof.

\subsection{The two small composite inputs}
For the first component use $p=5,43$, with one low digit in each
coordinate. Table~\ref{tab:215} specifies all 17 points and widths.
Every width is its displayed numerator divided by 10,000.
\begin{table}[ht]
\centering
\begin{tabular}{rl}\toprule
Width numerator&Low points $(a_5,a_{43})$\\\midrule
1181&$(0,18),(3,25),(1,27),(4,29)$\\
1181&$(0,24),(0,22),(3,42),(2,5)$\\
579&$(1,19),(2,40),(3,22),(4,35)$\\
1476&$(3,3),(4,39)$\\
984&$(0,33),(1,28),(2,20)$\\\bottomrule
\end{tabular}
\caption{The complete 17-point input at $5$ and $43$.}\label{tab:215}
\end{table}
For the edge relation in \eqref{eq:lowedge}, the graph has 78 directed
edges and maximum weighted path $9992/10000<1$. Thus Lemma~\ref{lem:paths}
provides the intervals. Its low-point moment is
\[
 Z_{215}(f)=8(1181/10000)^f+4(579/10000)^f
       +2(1476/10000)^f+3(984/10000)^f.
\]
With $f=0.0268$, exact arithmetic gives
\begin{equation}\label{eq:215moment}
 \log Z_{215}(f)-(2\al-1)\log215>0.00008383.
\end{equation}
By Lemma~\ref{lem:odd}, the resulting alphabet in base $215^2$ therefore
has moment greater than $(215^2)^{\al}$.

For the second component use $p=19,23$, with two low digits at each
prime. Its 729 points have a short description. Let
\[
 S=\{(a_i,b_i):0\le i<27\}\subseteq\mathbb F_{19}\times\mathbb F_{23}
\]
be the seed in Table~\ref{tab:affine}. The same table specifies four
integers $(A_i,B_i,C_i,D_i)$ for each seed point. Take all ordered pairs
$(i,j)$ and set
\begin{equation}\label{eq:437points}
 \begin{split}
 (a_{19,0},a_{19,1})&=(a_i,A_i a_j+B_i\bmod19),\\
 (a_{23,0},a_{23,1})&=(b_i,C_i b_j+D_i\bmod23).
 \end{split}
\end{equation}
Thus each first low digit selects an affine copy of the same seed for
the second low digit. The multipliers $A_i,C_i$ are nonzero squares in
their respective fields. They preserve the square relation inside a
fiber. Varying the copies between fibers changes the constraints imposed
by pairs whose prime coordinates first differ at different positions.
The proof uses the explicit points \eqref{eq:437points}; no general
claim that every such choice is feasible is needed.

The accompanying array assigns each point a positive width $n_{ij}/10000$.
These widths are chosen jointly, not by multiplying seed widths.
The full graph has 109,116 edges; its maximum weighted path is
$9972/10000<1$. For $f=0.0673$ the exact moment check gives
\begin{equation}\label{eq:437moment}
 \log\sum_{i,j=0}^{26}(n_{ij}/10000)^f
       -2(2\al-1)\log437>0.00025511.
\end{equation}
Lemma~\ref{lem:odd} now gives an alphabet in base $437^4$ with moment
greater than $(437^4)^{\al}$. The table of 729 width numerators has
165 distinct values and is supplied as machine-readable data rather
than as several pages of integers. Section~\ref{sec:verification}
specifies its format and the exact checks that certify both statements.


\section{A binary recursion with parity windows}\label{sec:binary}
The prime 2 does not furnish an odd-prime chain. Its useful local fact is
instead that every odd square is 1 modulo 8. A branch fixes a digit modulo
4, and the parity of its remaining tail determines the digit modulo 8.
Thus two windows per state contain exactly the information needed to order
square differences between branches.

\subsection{States and transitions}
A state $s$ has an even window $E_s$ and an odd window $O_s$ in $[0,1]$.
A window may be absent; each present window is a closed interval of
positive length, and at least one window is present. At depth $m$, the
state will describe an ordered interval alphabet modulo $4^m$, with the
interval of each digit inside the window of its parity.

A transition from $s$ has at most one branch for each $r\in\{0,1,2,3\}$.
A branch specifies a child $j$, scale $u>0$, shift $t\in\R$, and bits
$\tau,\eta\in\{0,1\}$, called swap and reflect. For a child digit $x$
modulo $Q=4^{m-1}$, let
\begin{equation}\label{eq:binarydigit}
 x'=((-1)^\eta x+\tau)\bmod Q,\qquad y=r+4x'.
\end{equation}
When $\eta=1$ reflect the child interval $[a,a+w]$ to
$[1-a-w,1-a]$; then apply the increasing affine map $z\mapsto t+uz$.
The parity of $x'$ is the parity of $x$ toggled by $\tau$, since $Q$ is
even. For transformed parity $p\in\{0,1\}$, let $J_{r+4p}$ be the
placed window obtained from the child window of parity $p\mathbin{\mathrm{xor}}\tau$,
with the same reflection and affine map. An absent child window contributes
no $J_{r+4p}$.

The transition conditions are
\begin{align}
 J_{r+4p}&\subseteq E_s\quad(r\text{ even}),\qquad
 J_{r+4p}\subseteq O_s\quad(r\text{ odd}),\label{eq:containment}\\
 \sup J_k&\le\inf J_{k+1\bmod8}
 \quad\text{whenever both windows are present}.\label{eq:cyclic}
\end{align}
In particular, \eqref{eq:cyclic} includes the comparison from class $7$
to class $0$. Missing classes make such a cyclic specification feasible.

\begin{lemma}[Binary transition]\label{lem:binary}
For $m\ge2$, if the child states are valid at depth $m-1$ and
\eqref{eq:containment}--\eqref{eq:cyclic} hold, the branches construct a
valid state at depth $m$.
\end{lemma}
\begin{proof}
Different branches give different residues modulo four, and within a
branch \eqref{eq:binarydigit} is injective. Thus the digits are distinct.
Containment of intervals follows from \eqref{eq:containment}.

Suppose $Y-X$ is a nonzero square modulo $4^m$. If the digits belong to
one branch, then $Y-X=4(y'-x')$ modulo $4^m$. Any corresponding square root
is even, so division by four gives a square difference modulo $Q$.
Without reflection the child order applies directly. With reflection the
transformed difference equals the negative of the original child
difference; consequently the child square edge has the opposite orientation.
Reflection of the intervals reverses that order once more, giving the
required transformed interval order. Translation cancels in differences,
and the positive scale preserves order.

For different branches whose residues differ by two modulo four, the
difference is $2$ modulo four and is not a square. In the remaining case
the difference is odd. Every odd square is $1$ modulo eight, since
$(2k+1)^2=4k(k+1)+1$. Therefore the numerical classes of $X,Y$ are
$k,k+1$ cyclically modulo eight. Their intervals lie in the corresponding
$J$ windows, so \eqref{eq:cyclic} orders them. This exhausts all cases.
Only the necessary condition on odd squares is used.
\end{proof}

\subsection{Moment growth and a finite depth}
Fix $0<f<1$. Let $v_s>0$ and $\lambda>1$ satisfy
\begin{equation}\label{eq:rows}
 \sum_{\text{branches }s\to j}u^f v_j\ge \lambda v_s\quad\text{for every }s.
\end{equation}
This positive-vector certificate is a familiar device for bounding the
growth of a nonnegative matrix \cite{Seneta2006}. Here the required
consequence is just the following elementary induction.

\begin{lemma}\label{lem:growth}
Let $d_s$ be the length of the wider present window of state $s$, and put
$c=\min_s d_s/v_s>0$. The transitions construct depth-$m$ states with
moments $Z_{s,m}\ge c\lambda^{m-1}v_s$ for all $m\ge1$.
\end{lemma}
\begin{proof}
At depth one choose just the digit $0$ or $1$ of the parity of a widest
window (either choice is allowed in a tie) and assign it the whole window. A singleton has no nontrivial
square edge. Its moment $d_s^f$ is at least $d_s$, since $0<d_s\le1$.
This gives $Z_{s,1}\ge cv_s$. Reflection and translation do not change
widths, scaling multiplies each moment by $u^f$, and different branches
are disjoint. Thus
\[
 Z_{s,m}=\sum_{s\to j}u^fZ_{j,m-1}
 \ge c\lambda^{m-2}\sum_{s\to j}u^fv_j\ge c\lambda^{m-1}v_s.
\]
Validity follows inductively from Lemma~\ref{lem:binary}.
\end{proof}

If a root state has both windows $[0,1]$ and $v_{\mathrm{root}}=1$, it
produces an unrestricted interval alphabet modulo $4^m$.
Shrinking every final interval at its left endpoint by a factor of $1/2$
preserves all orders, makes all widths at most $1/2$, and multiplies the
moment by $2^{-f}$. Therefore the condition
\begin{equation}\label{eq:binarymargin}
 \log c+(m-1)\log \lambda-f\log2-\alpha m\log4>0
\end{equation}
implies the required moment $Z_f>(4^m)^\alpha$. This uses only a fixed
finite depth, even if its numerical value is enormous.

\subsection{The explicit seven-state example}
Use states $0,\ldots,6$, with root 1 and windows in
Table~\ref{tab:windows}. All 26 transitions are in
Table~\ref{tab:transitions}. Together these two tables specify the entire
recursion, including all shifts and reflections.
\begin{table}[ht]
\centering
\begin{tabular}{rll}\toprule
State&Even window $E_s$&Odd window $O_s$\\\midrule
0&$[0,1]$&absent\\
1&$[0,1]$&$[0,1]$\\
2&$[0,.22]$&$[.51,1]$\\
3&$[0,1]$&$[.22,.51]$\\
4&$[0,.33]$&$[.40,1]$\\
5&$[0,.52]$&$[.22,1]$\\
6&$[0,1]$&$[.33,.40]$\\\bottomrule
\end{tabular}
\caption{The seven parity-window states.}\label{tab:windows}
\end{table}

The root illustrates how the geometry works. Its nonempty placed windows
are
\[
\begin{array}{lll}
 J_0=[0,49/178],&J_1=[49/178,39/89],&J_2=[39/89,65/89],\\
 J_3=[65/89,1],&J_5=[0,50/89],&J_6=[50/89,1].
\end{array}
\]
Classes 4 and 7 are absent. The required comparisons are precisely
$J_0\le J_1\le J_2\le J_3$ and $J_5\le J_6$, where an inequality
means that the left interval ends before the right one starts.
They hold by the displayed endpoints. The other states are checked in
the same way. In total there are 86 endpoint containment and order
comparisons, all between short rational numbers.

Fix $f_B=0.154$ and write $\phi(t)=t^{f_B}$. The moment transition matrix,
with rows and columns ordered from 0 to 6, is explicitly
\[
 M=\begin{pmatrix}
 0&2&0&0&0&0&0\\
 u&0&0&\beta&0&\beta&0\\
 v&0&0&0&0&0&0\\
 w&0&1&0&0&0&1\\
 x&0&0&0&0&0&0\\
 y&z&0&0&0&0&0\\
 t&0&0&0&1&0&0
 \end{pmatrix}.
\]
Here
\[
\begin{aligned}
 u&=\phi(49/178)+\phi(24/89),&\beta&=\phi(50/89),\\
 v&=2\phi(.22)+2\phi(.49),&w&=2\phi(.11),\\
 x&=2\phi(.33)+2\phi(.60),&y&=\phi(.22)+\phi(.48),\\
 z&=2\phi(.15),&t&=1+2\phi(.07).
\end{aligned}
\]
To prove growth, we do not need to compute an eigenvalue. Use the exact
positive vector
\begin{equation}\label{eq:vector}
 V=(.69921,1,.82524,.94106,.86406,.93399,.87113)^{\mathsf T}
\end{equation}
and the rational number $\lambda=2.8602$. Exact bounds for the fractional
powers establish, in each of the seven rows,
\begin{equation}\label{eq:rowmargin}
 (MV)_s-\lambda V_s>0.0001195.
\end{equation}
In the growth lemma the initialization constant is
$c=12250/20631$, attained at state 2. Set $m=100000$. The two useful
numerical inequalities are
\begin{align}
 \log\lambda-\al\log4&>0.00007904,\label{eq:growthmargin}\\
 \log c+(m-1)\log\lambda-f_B\log2-\al m\log4&>6.225.
 \label{eq:depthmargin}
\end{align}
Consequently, halving the final widths produces a valid alphabet in
base $4^{100000}$ with moment greater than $(4^{100000})^{\al}$.
The depth is large but fixed; induction represents the alphabet, so its
digits never have to be enumerated. Every resulting width is positive
and at most $1/2$.


\section{Combining the nine components}\label{sec:parameters}
The six prime chains and their exact rational powers are in
Table~\ref{tab:chains}. For a chain of length $t$ modulo $p$,
Lemma~\ref{lem:chain} gives moment $pt^{1-f_p}$ in base $p^2$.
The required inequality follows from
\[
 (1-f_p)\log t-(2\al-1)\log p>0.
\]
Every row passes this exact comparison; these powers are deliberately
rounded down and are not individually optimized. Their sum is $0.51011$.
\begin{table}[ht]
\centering
\begin{tabular}{rrll}\toprule
$p$&$t$&Ordered chain&$f_p$\\\midrule
3&2&$0,1$&.18215\\
7&3&$0,1,2$&.08603\\
11&4&$0,1,4,5$&.10746\\
31&7&$0,1,8,5,2,9,10$&.08940\\
59&9&$0,1,28,17,22,4,26,20,29$&.04242\\
103&11&$0,1,29,92,2,61,17,30,93,18,19$&.00265\\\bottomrule
\end{tabular}
\caption{Six chains of square differences. All forward differences are
nonzero squares modulo $p$, and all backward differences are nonsquares.}
\label{tab:chains}
\end{table}

\begin{proof}[Proof of Theorem~\ref{thm:main}]
The component bases and contributions are
\[
\begin{array}{c|c}
\text{square bases}&\text{contributions}\\\hline
3^2,7^2,11^2,31^2,59^2,103^2&\text{the six powers in Table~\ref{tab:chains}}\\
215^2&.0268\\
437^4&.0673\\
4^{100000}&.154.
\end{array}
\]
Their underlying prime factors are respectively
$3,7,11,31,59,103,\{5,43\},\{19,23\},2$, so the bases are pairwise
coprime perfect squares. The chain lemma,
\eqref{eq:215moment}, \eqref{eq:437moment}, and
\eqref{eq:depthmargin} establish the required moment inequalities.
Their contributions add to $0.75821$, leaving the exact surplus
\[
 \sum_i f_i-\al=0.000209>0.
\]
Every alphabet is now finite and fixed. The chain widths are less than
one, both odd inputs have widths strictly between zero and one, and the
binary widths have been halved. Thus common constants
$0<\sigma\le\rho<1$ exist. Theorem~\ref{thm:criterion} applies and
proves the result for every sufficiently large $N$.
\end{proof}

The fixed binary depth is chosen before taking $N\to\infty$. Its size
affects constants and the threshold for $N$, but causes no uniformity
issue in the interval-moment criterion. In fact that criterion constructs
one fixed block whose repetitions give $D(N)\ge cN^{\al}$ for a fixed
$c>0$ and all sufficiently large $N$.

\section{What the finite certificate checks}\label{sec:verification}
The general proof has reduced the theorem to finite statements about
explicit integer arrays and rational numbers. This section explains their
verification. The attached Python program contains the complete data and
checkers, uses only the standard library, and makes no network requests.

\paragraph{Odd geometry.}
The files \path{data/q215.json} and \path{data/q437.json}
list the low points and their width numerators. For a list of primes,
each point contains one low word for each prime, in increasing digit
significance. The common width denominator is 10,000. Thus the latter
file specifies all $n_{ij}$ in \eqref{eq:437moment}, with the point itself
identifying its pair $(i,j)$. A separate seed-and-map check regenerates
\eqref{eq:437points} and verifies that it agrees with the supplied array.

The checker verifies primality, point distinctness, digit ranges, positive
widths strictly below one, and the full edge relation \eqref{eq:lowedge}
on every ordered pair of distinct points. All edges must point forward
in the supplied ordering. For width numerator $n_j$, compute
\[
 d_j=n_j+\max\bigl(\{d_i:i\rightsquigarrow j\}\cup\{0\}\bigr).
\]
Induction in that ordering shows that $d_j$ is the maximum path numerator
ending at $j$. The check $\max_jd_j\le10000$ therefore proves feasibility
by Lemma~\ref{lem:paths}, and also defines the interval starts as
$(d_j-n_j)/10000$. This checks all 78 and 109,116 edges; no graph sampling
or assumption about a transitive reduction is involved.

\paragraph{Binary geometry and prime chains.}
The file \path{data/binary.json} records the windows,
branches, vector, growth constant, moment power and depth displayed above.
The checker constructs every placed window $J_k$ with exact fractions and
checks all conditions \eqref{eq:containment}--\eqref{eq:cyclic}. It also
checks the valid state indices, unique branch residues, positive scales,
and the two bits. These 86 comparisons, together with the binary lemma,
certify every depth. A direct check at depth 5 is an additional sanity
check and is not used in place of induction. The six chains are checked
in both orientations on all 244 ordered pairs of distinct entries.
Primality and disjointness of all CRT prime factors are checked as well.

\paragraph{Exact arithmetic for the moments.}
For completeness we give the elementary enclosure used to check fractional
powers. If $1\le x\le2$ and $z=(x-1)/(x+1)$, then $0\le z\le1/3$ and
\begin{equation}\label{eq:logseries}
 L_n(x):=2\sum_{j=0}^{n-1}\frac{z^{2j+1}}{2j+1}
 \le\log x\le
 L_n(x)+\frac{2z^{2n+1}}{(2n+1)(1-z^2)}.
\end{equation}
Indeed, integrate the geometric series to expand
$\log((1+z)/(1-z))$. Bound the tail by replacing its denominators by
$2n+1$ and summing a geometric series. For an arbitrary positive rational
$x$, write $x=2^k y$, $1\le y\le2$, and bound
$k\log2+\log y$, reversing endpoints when multiplying by a negative $k$.
All these bounds are rational.

Write $\ell(x),u(x)$ for the resulting lower and upper bounds on
$\log x$. A proposed positive rational lower bound $L$ for $x^f$, $f\ge0$,
is accepted only if
\[
 u(L)\le f\ell(x).
\]
This proves $L\le x^f$ by monotonicity of the logarithm.
The odd checker uses 32 terms of \eqref{eq:logseries} and the binary
checker uses 70. Decimal arithmetic proposes enclosures; exact rational
comparisons accept or reject them. Thus the validity of a successful check
does not depend on the accuracy of those proposals.

Summing accepted lower bounds for the odd powers and using upper bounds
for the subtracted logarithms proves \eqref{eq:215moment} and
\eqref{eq:437moment}. The same procedure proves every row of
\eqref{eq:rowmargin} and the fixed-depth inequality
\eqref{eq:depthmargin}. Each chain power is checked directly as a rational,
and the final contribution surplus is exact rational addition.
No spectral-radius approximation or asymptotic estimate enters these
acceptance tests.

\paragraph{Scope.}
The computer-assisted step is the execution of these finite integer and
rational checks. The mathematical reasons that their acceptance implies
the theorem are proved in the preceding sections. The certificate does
not assert that the inputs are optimal. Its complete source and data can
be extracted and inspected; no missing research notes or search programs
are premises of the argument.

\appendix
\section{Explicit local tables}\label{app:tables}
Table~\ref{tab:affine} supplies all seed points and conditional maps in
\eqref{eq:437points}. Each row $i$ specifies the first low point
$(a_i,b_i)$ and the map
$(a,b)\mapsto(A_i a+B_i,C_i b+D_i)$ in
$\mathbb F_{19}\times\mathbb F_{23}$ for its second low point.
Table~\ref{tab:transitions} uses the branch convention
\eqref{eq:binarydigit}. A blank branch is absent; every listed decimal
is exact.

\begin{table}[htbp]
\centering\small
\begin{tabular}{rrrrrrr}\toprule
$i$&$a_i$&$b_i$&$A_i$&$B_i$&$C_i$&$D_i$\\\midrule
0&10&6&1&0&1&0\\
1&14&6&1&0&12&7\\
2&2&6&1&0&13&11\\
3&10&19&1&0&1&0\\
4&0&16&1&0&1&0\\
5&0&19&7&6&12&8\\
6&11&16&1&0&13&10\\
7&0&20&7&6&1&0\\
8&11&19&1&0&1&0\\
9&11&20&1&0&2&10\\
10&10&9&1&0&1&0\\
11&16&10&1&0&1&0\\
12&16&19&1&0&16&17\\
13&10&22&1&0&1&0\\
14&17&10&1&0&2&10\\
15&16&22&1&0&8&20\\
16&17&19&1&0&2&7\\
17&14&9&1&9&4&19\\
18&17&22&5&1&3&17\\
19&14&22&1&9&1&0\\
20&2&9&5&2&8&1\\
21&2&22&1&0&6&20\\
22&16&12&1&0&1&0\\
23&0&5&1&0&1&0\\
24&1&12&1&0&16&11\\
25&11&5&1&0&13&10\\
26&17&12&1&0&1&0\\
\bottomrule\end{tabular}
\caption{The 27-point seed and its conditional affine maps for the
729-point component.}\label{tab:affine}
\end{table}

\begin{table}[htbp]
\centering\small
\renewcommand{\arraystretch}{1.13}
\begin{tabular}{rrrllrr}\toprule
Parent $s$&$r$&Child $j$&Scale $u$&Shift $t$&Swap $\tau$&Reflect $\eta$\\\midrule
0&0&1&$1$&$0$&0&0\\
0&2&1&$1$&$0$&0&0\\
\midrule
1&0&0&$\frac{49}{178}$&$0$&0&0\\
1&1&3&$\frac{50}{89}$&$0$&1&1\\
1&2&5&$\frac{50}{89}$&$\frac{39}{89}$&0&0\\
1&3&0&$\frac{24}{89}$&$\frac{65}{89}$&0&0\\
\midrule
2&0&0&$0.22$&$0$&0&0\\
2&1&0&$0.49$&$0.51$&1&0\\
2&2&0&$0.22$&$0$&0&0\\
2&3&0&$0.49$&$0.51$&0&0\\
\midrule
3&0&2&$1$&$0$&0&0\\
3&1&0&$0.11$&$0.22$&0&0\\
3&2&6&$1$&$0$&1&0\\
3&3&0&$0.11$&$0.4$&0&0\\
\midrule
4&0&0&$0.33$&$0$&0&0\\
4&1&0&$0.6$&$0.4$&1&0\\
4&2&0&$0.33$&$0$&0&0\\
4&3&0&$0.6$&$0.4$&0&0\\
\midrule
5&0&0&$0.22$&$0$&0&0\\
5&1&1&$0.15$&$0.22$&0&0\\
5&2&1&$0.15$&$0.37$&0&0\\
5&3&0&$0.48$&$0.52$&0&0\\
\midrule
6&0&4&$1$&$0$&0&0\\
6&1&0&$0.07$&$0.33$&0&0\\
6&2&0&$1$&$0$&1&0\\
6&3&0&$0.07$&$0.33$&0&0\\
\bottomrule\end{tabular}
\caption{All 26 branches of the seven-state recursion. Each row means
$x'=((-1)^\eta x+\tau)\bmod4^{m-1}$, $y=r+4x'$, with the interval
reflected if $\eta=1$ and then mapped by $z\mapsto t+uz$.}
\label{tab:transitions}
\end{table}

\clearpage
\section{Reproducing the checks}\label{app:reproduction}
The \LaTeX\ source contains the complete standalone verifier. Compilation
writes it to a separate Python file and attaches the program to the PDF:
\begin{center}
\attachfile[description={Exact finite data and standard-library Python verifier},
 mimetype=text/plain,icon=Paperclip]{sarkozy-simple-certificate.py}
\quad\textbf{Complete certificate and verifier}
\end{center}
Save the attachment, or use the separate copy supplied with the paper,
and run
\begin{lstlisting}[language=bash]
python3 sarkozy-simple-certificate.py
\end{lstlisting}
To retain the unpacked source, JSON inputs and reports in a fresh directory,
run
\begin{lstlisting}[language=bash]
python3 sarkozy-simple-certificate.py --extract certificate-check
\end{lstlisting}
The program rejects Python's \texttt{-O} and \texttt{-OO} modes, so all
verification assertions execute. It checks the identities of the bundled
inputs as well as their mathematical conditions. It needs standard Python 3
and no numerical optimization packages.

To compile from the single \LaTeX\ source, use an ordinary TeX Live
installation and run the following command until cross-references settle
(three runs suffice):
\begin{lstlisting}[language=bash]
pdflatex -interaction=nonstopmode -halt-on-error \
  -no-shell-escape sarkozy_simple_0.758.tex
\end{lstlisting}
The bibliography and certificate are embedded in that source; no external
bibliography or certificate download is needed. The tables and explicit
parameters above are part of the mathematical specification, and the
attachment provides the remaining width array and executable checks.

\begingroup
\raggedright
\begin{thebibliography}{99}

\bibitem{Furstenberg1977}
H. Furstenberg,
\emph{Ergodic behavior of diagonal measures and a theorem of
Szemer\'edi on arithmetic progressions},
J. Analyse Math. \textbf{31} (1977), 204--256.
\href{https://doi.org/10.1007/BF02813304}{doi:10.1007/BF02813304}.

\bibitem{Sarkozy1978}
A. S\'ark\"ozy,
\emph{On difference sets of sequences of integers. I},
Acta Math. Acad. Sci. Hungar. \textbf{31} (1978), no.~1--2, 125--149.
\href{https://doi.org/10.1007/BF01896079}{doi:10.1007/BF01896079}.

\bibitem{Ruzsa1984}
I.~Z. Ruzsa, \emph{Difference sets without squares},
Period. Math. Hungar. \textbf{15} (1984), no.~3, 205--209.
\href{https://doi.org/10.1007/BF02454169}{doi:10.1007/BF02454169}.

\bibitem{PintzSteigerSzemeredi1988}
J. Pintz, W.~L. Steiger, and E. Szemer\'edi,
\emph{On sets of natural numbers whose difference set contains no squares},
J. London Math. Soc. (2) \textbf{37} (1988), no.~2, 219--231.
\href{https://doi.org/10.1112/jlms/s2-37.2.219}{doi:10.1112/jlms/s2-37.2.219}.

\bibitem{BeigelGasarch2008}
R. Beigel and W. Gasarch,
\emph{Square-difference-free sets of size $\Omega(n^{0.7334\ldots})$},
2008, \href{https://arxiv.org/abs/0804.4892v3}{arXiv:0804.4892v3}.

\bibitem{Lewko2015}
M. Lewko,
\emph{An improved lower bound related to the
Furstenberg--S\'ark\"ozy theorem},
Electron. J. Combin. \textbf{22} (2015), no.~1, Paper~1.32, 6~pp.
\href{https://doi.org/10.37236/4656}{doi:10.37236/4656}.

\bibitem{BloomMaynard2022}
T.~F. Bloom and J. Maynard,
\emph{A new upper bound for sets with no square differences},
Compos. Math. \textbf{158} (2022), no.~8, 1777--1798.
\href{https://doi.org/10.1112/S0010437X22007679}{doi:10.1112/S0010437X22007679}.

\bibitem{GreenSawhney2025}
B. Green and M. Sawhney,
\emph{New bounds for the Furstenberg--S\'ark\"ozy theorem},
2025, \href{https://arxiv.org/abs/2411.17448v2}{arXiv:2411.17448v2}
(first version 2024).

\bibitem{Krachun2026}
D. Krachun,
\emph{Square-difference-free sets beyond the three-quarter barrier},
2026, \href{https://arxiv.org/abs/2608.01325v1}{arXiv:2608.01325v1}.

\bibitem{Stanley1986}
R.~P. Stanley,
\emph{Two poset polytopes},
Discrete Comput. Geom. \textbf{1} (1986), 9--23.
\href{https://doi.org/10.1007/BF02187680}{doi:10.1007/BF02187680}.


\bibitem{Seneta2006}
E. Seneta,
\emph{Non-negative Matrices and Markov Chains}, second ed.,
Springer Series in Statistics, Springer, New York, 1981,
reprinted 2006.
\href{https://doi.org/10.1007/0-387-32792-4}{doi:10.1007/0-387-32792-4}.


\bibitem{NaslundCompanion2026}
E.~Naslund,
\emph{Square-difference-free sets of exponent $0.75806746$: ordered
intervals, composite digit codes, and a binary recursion},
unpublished companion manuscript, September 2026.

\end{thebibliography}

\endgroup

\end{document}
